\documentclass[11pt,a4paper]{article}
\usepackage[T1]{fontenc}
\usepackage{lmodern,amsmath,amssymb,amsthm,booktabs,array,microtype}
\usepackage[margin=26mm,top=24mm,bottom=25mm,headheight=14pt]{geometry}
\usepackage{xcolor,fancyhdr,titlesec,enumitem}
\definecolor{ink}{HTML}{18354A}
\definecolor{accent}{HTML}{14665F}
\usepackage[colorlinks=true,linkcolor=accent,citecolor=accent,urlcolor=accent,pdftitle={Daisy chains and strong chaplets},pdfauthor={Research manuscript prepared for Keio}]{hyperref}
\urlstyle{same}
\titleformat{\section}{\large\bfseries\color{ink}}{\thesection}{.7em}{}
\titleformat{\subsection}{\normalsize\bfseries\color{ink}}{\thesubsection}{.7em}{}
\pagestyle{fancy}\fancyhf{}
\lhead{\small\color{ink} Daisy chains and strong chaplets}
\rhead{\small\color{ink} Research manuscript}
\cfoot{\small\thepage}
\newtheorem{theorem}{Theorem}[section]
\newtheorem{lemma}[theorem]{Lemma}
\newtheorem{proposition}[theorem]{Proposition}
\newtheorem{corollary}[theorem]{Corollary}
\newcommand{\U}{U}
\newcommand{\F}{\mathbb F}
\newcommand{\Z}{\mathbb Z}
\newcommand{\id}{\mathrm{id}}
\newcommand{\st}{\,\colon\,}
\setlength{\parindent}{1em}
\setlength{\parskip}{.22em}
\raggedbottom
\begin{document}
\thispagestyle{plain}
\noindent{\sffamily\small\bfseries\color{accent} ADDITIVE COMBINATORICS}\hfill
{\small 21 November 2026}\par
\vspace{.7em}
\noindent{\LARGE\bfseries\color{ink} Daisy chains and strong chaplets}\par
\vspace{.25em}
\noindent{\large\color{ink} Constructions for Sun's and Preece's existence problems}\par
\vspace{.6em}
\noindent{\small Research manuscript prepared for Keio}\par
\vspace{.8em}
\noindent\textbf{Abstract.}
We construct a cyclic ordering of the units modulo every odd integer $n>1$ in which the adjacent differences run through the same units exactly once. This proves the conjecture in Remark~1.4 of Zhi-Wei Sun's \emph{Some new problems in additive combinatorics} for its full stated range. Preece called this ordering a daisy chain. Using signed representatives, the same method constructs a strong chaplet for every odd modulus greater than five other than a power of three. This answers Preece's 2009 existence question. Explicit robust chaplets at 193 and 257 also answer his question about a prime of the form $2^i+1$. We also determine exactly when a chain and its first two difference sequences can all list the units: the unique exceptional odd modulus is $15$. The construction prescribes multiplication along the fibres of a cyclic ordering. An explicit change at one edge makes the return map transitive while all edge labels remain bijective on their fibres. We give an optimal simultaneous version for several linear forms. For a finite commutative ring of odd order, we also determine the minimum number of cycles in a permutation of the units whose differences permute the units.
\section{The original problem}
For $n>1$, write $\U(n)=(\Z/n\Z)^\times$ and $L=\varphi(n)$. A word $a=(a_i)_{i\in\Z/L\Z}$ has forward differences $(\Delta a)_i=a_{i+1}-a_i$, including the closing edge. Write $[a]$ for the multiset of entries. Preece called $a$ a \emph{daisy chain} if $[a]=[\Delta a]=\U(n)$, with each unit counted once \cite{preece08}. Sun used the opposite difference. Negation permutes the units, so the two formulations are equivalent for every cyclic word.
\begin{theorem}[Sun's conjecture]\label{thm:sun}
For every odd integer $n>1$, there is a cyclic word $a$ with
\[
                 [a]=[\Delta a]=\U(n).
\]
\end{theorem}
Sun dated the conjecture 4 October 2013 in his lecture notes \cite{sun-talk} and published it in Remark~1.4 following his odd-prime-power theorem \cite{sun}. Prime powers admit a particularly strong construction. If $g$ generates $\U(p^b)$ for an odd prime $p$, then $a_i=g^i$ has $\Delta^s a_i=(g-1)^s g^i$ for all $s\ge0$. Indeed $g\bmod p$ generates $\F_p^\times$, so $g-1$ is a unit. This is Preece's Theorem~2.1 and the observation opening his Section~6.
\par
The first composite obstruction arises already at $n=15$. The group $\U(15)$ has order eight and its elements have order at most four. Iterating a fixed multiplier cannot visit all eight units. The problem is to join local cyclic orders without repeating any difference.
\par
Preece's second problem uses half as many vertices. Can $\varphi(n)/2$ units be arranged in a cycle whose differences are distinct up to sign, for every odd $n>5$ other than a power of three? He emphasizes that the number of prime divisors is unbounded \cite[Section~10]{chaplet}. Theorem~\ref{thm:chaplet} answers yes: the entries and their differences can each represent all the opposite pairs of units. The further requirement that the two representative sets agree up to overall sign is called robustness. Section~\ref{sec:robust} supplies explicit certificates for the primes $193$ and $257$ singled out in the same paper.
\clearpage
\section{Prescribing the return map}
Consider an auxiliary construction problem. A cycle on $X$ has distinct difference labels. Replace each vertex by a copy of a cyclic multiplicative group $V$. Which changes of multiplier along the cycle join all copies in one orbit and preserve all labels? The following criterion gives both conditions separately.
\begin{lemma}[Fibre criterion]\label{lem:lift}
Let $x_0,\ldots,x_{L-1}$ list a subset $X$ of an abelian group, with $d_i=x_{i+1}-x_i$ listing a set $D$ once each. Let $V$ be the cyclic unit group of a finite commutative ring $S$. Put $r_i\in V$ and define
\[
                   T(x_i,y)=(x_{i+1},r_i y),\qquad y\in V .
\]
Then $T$ is a single cycle precisely when $R=\prod_{i=0}^{L-1}r_i$ generates $V$. Its difference labels list $D\times V$ precisely when every $r_i-1$ is a unit in $S$.
\end{lemma}
\begin{proof}
An orbit returns to the fibre over $x_i$ after exactly $L$ steps in the base. On that fibre $T^L$ multiplies by $R$, so each orbit has length $L\operatorname{ord}(R)$. All labels above $d_i$ have second coordinate $(r_i-1)y$. Multiplication by this coefficient permutes $V$ if it is a unit. If it is a nonunit, its value at $y=1$ is already outside $V$. The $d_i$ distinguish the fibres of labels.
\end{proof}
\begin{proof}[Proof of Theorem~\ref{thm:sun}]
Factor $n$ into powers of distinct primes and start with the smallest prime power. The powers of a primitive root give the initial cycle. At any later stage the old length $L$ is even and the next factor is $q=p^b$ with $p\ge5$. Choose a primitive root $g\bmod q$ and put
\[
              r_0=\cdots=r_{L-2}=-1,\qquad r_{L-1}=-g .
\]
Their product is $g$. The two possible difference factors are $-2$ and $-g-1$. The reduction of $g$ has order $p-1\ge4$, so both factors are units. Lemma~\ref{lem:lift} gives a single cycle with every difference represented once in the product. The Chinese remainder isomorphism preserves addition, multiplication and units, so the resulting word has the required property for the enlarged modulus. Induction includes every factor of $n$.
\end{proof}
\noindent\textbf{Example.} Starting from $(1,2)\bmod3$, use the two multipliers $-1,-2\bmod5$. Reading the resulting cycle and its forward differences modulo $15$ gives
\[
\begin{array}{c|rrrrrrrr}
a&1&14&7&8&4&11&13&2\\
\Delta a&13&8&1&11&7&2&4&14 .
\end{array}
\]
Every unit appears once in each row. Section~\ref{sec:two-differences} shows why no ordering modulo $15$ achieves this for a third row.
\section{Several linear forms on the same cycle}
For a unit $c$, write $D_c a_i=a_{i+1}-ca_i$. Thus $D_1$ is the difference and $D_{-1}$ is the sum. We call a set $C\subseteq\U(n)$ \emph{separated modulo the prime factors} if $c-c'$ is a unit for every distinct $c,c'\in C$. This is precisely injectivity of reduction on $C$ at each prime dividing $n$.
\begin{theorem}\label{thm:slopes}
If $n>1$ is odd and $\ell$ is its smallest prime factor, the largest size of such a set $C$ for which a single cyclic listing $a$ of $\U(n)$ satisfies
\[
                          [D_c a]=\U(n)\qquad(c\in C)
\]
is $\ell-2$. A maximizing set can contain $1$, and can contain both $1$ and $-1$ if $\ell\ge5$.
\end{theorem}
\begin{proof}
For any fixed $i$, the nonzero residue $a_{i+1}/a_i\bmod\ell$ is excluded from the image of $C$, since $D_ca_i$ must be a unit. Distinct members of $C$ have different nonzero reductions. This gives the upper bound $\ell-2$.
\par
Start with the factor $\ell^b$ and a generator $g_0$ of its units. Use the power cycle and the $\ell-2$ residues in $\F_\ell^\times\setminus\{\bar g_0\}$. On this cycle each allowed label is $(g_0-c)a_i$, a permutation of the units. For a later factor $q=p^b$ with generator $g$, use the exponent word
\[
                   -1,\ 2,\ \underbrace{1,-1,\ldots,1,-1}_{L-2}
\]
and multipliers $r_i=g^{e_i}$. The sum of its exponents is one. For this coordinate choose $\ell-2$ distinct nonzero residues outside $\{\bar g^{-1},\bar g,\bar g^2\}$ in $\F_p$. This is possible since $p-4\ge\ell-2$. Assign these residues bijectively to the initial list, lift them to units modulo $q$, and use CRT to obtain the global $c$'s. At the new stage the fibre coordinate of $D_c$ is $(r_i-c)y$, which permutes the units for each $i$. Its base coordinate already permutes all base labels. The return multiplier $g$ gives a single cycle for all choices of $c$ at once.
\par
The residue $1$ is allowed at every coordinate. If $\ell\ge5$, $-1$ is allowed initially, and $p\ge7$ thereafter implies that a generator and its inverse and square all differ from $-1$ modulo $p$. Keep $\pm1$ paired with themselves in all coordinates and choose their indicated lifts. This gives the final assertion.
\end{proof}
\begin{corollary}
For any integer $n>1$, the same cyclic ordering can have both adjacent sums and adjacent differences listing $\U(n)$ if and only if $\gcd(n,6)=1$.
\end{corollary}
\begin{proof}
Theorem~\ref{thm:slopes} gives sufficiency. For $2\mid n$ two units have even difference. For $3\mid n$ two units modulo three are either equal or negatives of each other, so their sum or their difference vanishes there.
\end{proof}
\section{A sharp classification for two successive differences}
\label{sec:two-differences}

For a cyclic word $a=(a_i)$ write $(\Delta a)_i=a_{i+1}-a_i$, with subscripts read cyclically. In Section~6 of his 2008 paper, Preece \cite{preece08} asks when the differences of a daisy chain themselves form a daisy chain. His term \emph{$t$-fertile} specifies the exact first point of failure. We say \emph{fertility at least two} when $a,\Delta a,\Delta^2a$ each list all the units once. Preece gives examples of several exact fertilities. The following existence classification uses the same iteration.

\begin{theorem}\label{thm:two-orders}
For every odd integer $n>1$, there is a cyclic listing of $\U(n)$ of fertility at least two if and only if $n\ne15$.
\end{theorem}

\begin{lemma}\label{lem:two-fibre}
Suppose $a$ is a cyclic listing of $\U(m)$ of fertility at least two and length $L$. Let $q=p^b$ be an odd prime power coprime to $m$. Choose a cyclic list $r_0,\ldots,r_{L-1}\in\U(q)$ such that their product generates $\U(q)$ and, for all $i$,
\[
 r_i-1\in\U(q),\qquad r_ir_{i+1}-2r_i+1\in\U(q).
\]
Then $\U(mq)$ admits a cyclic listing of fertility at least two.
\end{lemma}

\begin{proof}
On the pairs $(a_i,y)$ with $y\in\U(q)$, the successor map is $(a_{i+1},r_i y)$. After $L$ steps the second coordinate has been multiplied by the product of the $r_i$, so all pairs belong to one cycle. At a position $(a_i,y)$ the second coordinates of the first and second differences are $(r_i-1)y$ and $(r_i r_{i+1}-2r_i+1)y$, respectively. For fixed $i$ each ranges bijectively over the fibre. For each order $j=0,1,2$, the first coordinates $(\Delta^j a)_i$ already list the units modulo $m$. The Chinese remainder theorem gives the claim.
\end{proof}

If $p\ge7$ and $L\ge2$ is even, choose a primitive residue $g$ modulo $p$ outside $\{-3,-1/3\}$. For $p=7$ take $g=3$. For $p\ge11$ there are at least four primitive residues, since $\varphi(k)>2$ for every even $k\ge10$. A generator modulo $p$ has order at least six here, so it also differs from $-1$. Lift $g$ to a primitive root $G$ modulo $p^b$. Such a lift exists because a generator in the cyclic quotient of a cyclic group lifts to a generator. Use $r_i=-1$ everywhere except at one position, where $r_i=-G$. The product is $G$. On ordinary adjacent pairs the second coefficient is $4$. The two pairs containing the exceptional entry have second coefficients $3G+1$ and $G+3$. All required coefficients are units. This proves the lifting step also when $L=2$, when there are no ordinary adjacent pairs.

For $p=5$ and even $L\ge4$, use the following periodic list of residues modulo $5$, with exactly $L$ entries,
\[
           3,\ 4,\ 2,\ 3,\ \underbrace{2,\dots,2}_{L-4}.
\]
Among $\{2,3,4\}$ the only disallowed adjacent pairs for the second coefficient are $2\longrightarrow4$ and $4\longrightarrow3$. The displayed list avoids them, including across the join. All first coefficients are nonzero. Its product is $72\cdot2^{L-4}\pmod5$, a generator since $72\equiv2\pmod5$ and $L-4$ is even. To work modulo $5^b$, choose arbitrary unit lifts of all but one entry and choose the remaining entry to make the total product a primitive root lifting this residue. Its prescribed reduction is then automatic.

\begin{lemma}\label{lem:repeated-prime-difference}
Let $p\mid m$ be prime. If each of $a,\Delta a,\ldots,\Delta^h a$ lists $\U(m)$, then a cyclic word with the same property exists modulo $pm$.
\end{lemma}

\begin{proof}
Let $L=\varphi(m)$. Lift every $a_j$ arbitrarily to an integer $A_j$ and set $u=1+m$. Use the word
\[
        b_{j+tL}=u^t A_j\pmod{pm}
        \qquad(0\le j<L,\ 0\le t<p).
\]
The binomial theorem gives $u^t\equiv1+tm\pmod{pm}$, so $u$ has order $p$. For every unit modulo $pm$, its orbit under multiplication by $u$ is the entire fibre above its reduction modulo $m$. With cyclic indices one has $b_{k+L}=ub_k$. Subtraction preserves this identity at every order, including across block boundaries. For fixed $j$, the $p$ values of $(\Delta^s b)_{j+tL}$ reduce to $(\Delta^s a)_j$ and run through that unit's entire fibre. This proves the result for each $s\le h$.
\end{proof}

For $n=75$ we give a small certificate. Index the word by $0\le j+10t<40$, with $0\le j<10$. Its coordinates modulo $3$ and $25$ are
\[
  (-1)^{j+10t},\qquad 7^t b_j,\qquad
  b=(1,2,4,13,16,22,14,18,6,8).
\]
The group $H=\langle7\rangle=\{1,7,24,18\}\subset\U(25)$ has order four. For $s=0,1,2$, shifting by ten multiplies the second coordinate of $\Delta^s a$ by $7$. The table gives the least positive representatives of its $H$-cosets for the even and odd positions in the first block, in increasing order of position. The last two differences use $b_{10}=7$ and $b_{11}=14$.
\[
\begin{array}{c|cc}
 s & j=0,2,4,6,8 & j=1,3,5,7,9\\ \hline
 0 &(1,3,9,2,6)&(2,9,3,1,6)\\
 1 &(1,9,6,3,2)&(2,3,6,9,1)\\
 2 &(1,6,2,9,3)&(1,3,9,2,6)
\end{array}
\]
Each column at each order meets every coset. Modulo $3$ the entries of $\Delta^s a$ are $(-2)^s(-1)^j$, so their two parity classes have opposite residues. The table therefore certifies all three permutations modulo $75$.

\begin{proof}[Proof of Theorem~\ref{thm:two-orders}]
For a single odd prime power $q$, a cyclic list of successive powers of a primitive root $G$ has $s$th differences $(G-1)^s G^i$, which list all units for every $s$. If $3\nmid n$, begin with its smallest prime-power factor, then adjoin the others with Lemma~\ref{lem:two-fibre}. Each adjoined prime is at least seven. If $3\mid n$, begin with the full $3$-power factor, then adjoin all prime powers with prime at least seven, leaving the $5$-power for last. The current length before adjoining the latter is even and at least four except when the current modulus is $3$. In that remaining case $n=3\cdot5^b$. The certificate supplies $b=2$, and Lemma~\ref{lem:repeated-prime-difference} supplies each $b\ge3$.

It remains to exclude $n=15$. The entries of a daisy chain must alternate modulo $3$ because their first differences are units. For each of the two alternating layers the successor induces a permutation of $\mathbb F_5^\times$ on second coordinates. Call them $f$ and $h$. The two layers have different first differences modulo $3$, so $f(x)-x$ and $h(x)-x$ must also permute $\mathbb F_5^\times$.

We check that any such $f$ is a scalar. Extend by $f(0)=0$ and represent $f$ by its unique polynomial of degree at most four on $\mathbb F_5$. Both $f$ and $f-\mathrm{id}$ are permutations of the field. Thus $\sum_x f(x)=0$ and $\sum_x f(x)^2=\sum_x (f(x)-x)^2=\sum_x x^2=0$, whence $\sum_x x f(x)=0$. In $\mathbb F_5$ the sums of the first three positive powers vanish and the sum of fourth powers is $-1$. The first sum eliminates the quartic coefficient and the last eliminates the cubic coefficient. A quadratic of degree exactly two over a field of odd characteristic takes the same value at two points symmetric about its axis, so cannot be a permutation. The constant term vanishes, proving $f(x)=rx$. The same holds for $h(x)=sx$, and $r,s\in\{2,3,4\}$.

Two steps on a fixed layer multiply by $rs$. A full eight-cycle requires $rs$ to generate $\mathbb F_5^\times$. If $rs=2$, then $\{r,s\}=\{3,4\}$, and the second-difference coefficient $rs-2\cdot4+1$ vanishes. If $rs=3$, then $\{r,s\}=\{2,4\}$, and $rs-2\cdot2+1$ vanishes. In either case the second difference is zero modulo $5$ on one layer. This completes the exclusion.
\end{proof}
\section{Preece's existence question for chaplets}
\label{sec:chaplets}

Preece's \emph{Half-cycles and chaplets} \cite[Section~10, p.~279]{chaplet} asks about every odd $n>5$ other than powers of $3$. A chaplet is a cycle of $\varphi(n)/2$ distinct units whose consecutive differences are distinct and contain no opposite pair. In the constructions here the differences are units and contain one element from each pair $\{u,-u\}\subset\U(n)$. The chaplet is \emph{strong} if its entries also have this last property. We call such a choice of one representative from every pair a signed transversal.

\begin{theorem}\label{thm:chaplet}
Every odd integer $n>5$ that is not a power of $3$ admits a strong chaplet.
\end{theorem}

\begin{proof}
There are two lifting operations. First let $q\ge5$ be a prime not dividing the current modulus $m$, and let $a$ be a strong chaplet modulo $m$ of length $L$ whose differences are also a signed transversal of the units. Both properties will be maintained. Choose a primitive root $g$ modulo $q$. For odd $L$ use one multiplier $g$ and $(L-1)/2$ pairs $(g,g^{-1})$. For even $L$ use one pair $(g^{-1},g^2)$ and $(L-2)/2$ pairs $(g,g^{-1})$. In either case the product of the multipliers is $g$, and every multiplier $r_i$ differs from $1$, since $g$ has order at least four. Use the successor
\[
             (a_i,y)\ \longmapsto\ (a_{i+1},r_i y),
             \qquad y\in\U(q).
\]
The product condition gives one cycle. Its vertex set is $A\times\U(q)$ and its difference set is $D\times\U(q)$, where $A$ and $D$ are the sets of vertices and differences of the base cycle. Both products are signed transversals modulo $mq$. In fact, in each opposite pair $(v,w),(-v,-w)$, the choice is determined by which of $v,-v$ belongs to the base transversal.

Second, if $p\mid m$, the same block lift $b_{j+tL}=(1+m)^t A_j$ used in Lemma~\ref{lem:repeated-prime-difference} applies to signed transversals. Each base vertex and each base difference contributes exactly its full fibre of $p$ lifts. The full preimage of a signed transversal under reduction is again a signed transversal. Thus one may increase the exponent of any prime already present while preserving strength. For related special cases see Preece's Theorems~3.6 and~3.7 in the chaplet paper.

It remains to provide starting moduli.
\begin{itemize}
\item For prime $p\ge7$, the strong half-cycle of Buratti and Del Fra, reproduced as Theorem~2.1 in Preece's chaplet paper, is a strong chaplet modulo $p$, because every nonzero residue is a unit.
\item For a pure power $5^b$ with $b\ge2$, start at $25$ with Preece's strong chaplet (p.~256)
\[
                      (1,2,8,21,18,9,13,6,14,3).
\]
Its entries and differences have least absolute values $\{1,2,3,4,6,7,8,9,11,12\}$ modulo $25$.
\item If $3\mid n$, choose a prime $p\ge5$ dividing $n$ and a generator $g$ modulo $p$. In $\U(3)\times\U(p)$ put $x=(-1,g)$ if $p\equiv1\pmod4$, and $x=(-1,g^2)$ if $p\equiv3\pmod4$. In the first case $x$ has order $p-1$. An exponent with $g^k=-1$ is even, so $(-1,-1)$ is outside $\langle x\rangle$. In the second case $(p-1)/2$ is odd, hence $x$ has order $p-1$ and $-1$ is outside $\langle g^2\rangle$. In both cases $H=\langle x\rangle$ is an index-two subgroup avoiding the global minus identity. Also $x-1$ is a unit, because in the second case $p\ge7$. The cycle $(1,x,\dots,x^{p-2})$ modulo $3p$ has vertex set $H$ and difference set $(x-1)H$. Both are signed transversals. This is the index-two construction of Preece's Theorem~4.1 specialized to the required seed.
\end{itemize}
If $3\nmid n$ and $n$ has a prime divisor at least seven, use that prime as the starting modulus. The only other case with $3\nmid n$ is the pure $5$-power above. When $3\mid n$, start at $3p$ as constructed. All remaining distinct prime factors are at least five and can be adjoined by the first operation. Increasing the exponents by the second operation reaches $n$.
\end{proof}
\section{Robust chaplets and the exceptional small moduli}
\label{sec:robust}
A strong chaplet is \emph{robust} if its vertex set equals its difference set or its negative. In Section~10 of his 2009 paper \cite{chaplet}, Preece asks whether a robust chaplet exists for any prime of the form $2^i+1$, and explicitly discusses the case $193=3\cdot 2^6+1$. The cycle at $257=2^8+1$ below answers the existential question. Martin and Ollis later proved nonexistence at $17$ and reported a restricted search at $193$ and $257$ \cite[Section~6]{mo}. The two witnesses both realize equality with the forward difference set.

\begin{theorem}\label{thm:robust-certs}
There is a robust chaplet for each of the primes $193$ and $257$.
\end{theorem}

\begin{proof}
Read each list from left to right and from top to bottom, closing the final entry to the first. For modulus $p$, put $V=\{a_i\}$ and $D=\{a_{i+1}-a_i\bmod p\}$. Direct reduction of the following lists gives $D=V$ and $V\cup(-V)=\mathbb F_p^\times$. The numbers of distinct values in $V$ and $D$ are both $96$ for $p=193$, and both $128$ for $p=257$. These equalities verify each cyclic edge and both signed-transversal requirements.

\par\noindent\begin{minipage}{\linewidth}
\noindent\textbf{Modulus $193$.}\par\smallskip
{\small\setlength{\tabcolsep}{3pt}\centering
\begin{tabular}{*{16}{r}}
1 & 52 & 84 & 3 & 24 & 60 & 54 & 102 & 36 & 96 & 129 & 149 & 130 & 115 & 132 & 158 \\
182 & 38 & 93 & 10 & 8 & 49 & 103 & 99 & 55 & 126 & 178 & 188 & 124 & 30 & 151 & 191 \\
116 & 39 & 73 & 112 & 105 & 106 & 28 & 17 & 101 & 187 & 117 & 41 & 79 & 33 & 135 & 119 \\
22 & 40 & 85 & 164 & 95 & 32 & 137 & 26 & 184 & 59 & 45 & 140 & 50 & 23 & 31 & 104 \\
147 & 118 & 168 & 131 & 51 & 82 & 21 & 127 & 186 & 128 & 86 & 12 & 113 & 13 & 43 & 71 \\
156 & 20 & 166 & 110 & 123 & 34 & 57 & 48 & 179 & 174 & 177 & 189 & 18 & 146 & 121 & 68 \\
\end{tabular}\par}
\end{minipage}\par

\par\noindent\begin{minipage}{\linewidth}
\noindent\textbf{Modulus $257$.}\par\smallskip
{\small\setlength{\tabcolsep}{3pt}\centering
\begin{tabular}{*{16}{r}}
256 & 241 & 252 & 21 & 69 & 221 & 34 & 73 & 255 & 42 & 182 & 120 & 206 & 183 & 39 & 195 \\
229 & 32 & 63 & 103 & 99 & 87 & 203 & 179 & 133 & 227 & 125 & 186 & 204 & 169 & 245 & 189 \\
248 & 219 & 60 & 143 & 228 & 200 & 6 & 48 & 127 & 118 & 254 & 100 & 59 & 84 & 153 & 230 \\
222 & 156 & 175 & 11 & 140 & 135 & 31 & 199 & 85 & 106 & 52 & 25 & 224 & 19 & 109 & 129 \\
116 & 45 & 136 & 98 & 145 & 77 & 126 & 93 & 91 & 26 & 40 & 185 & 235 & 86 & 50 & 67 \\
177 & 47 & 92 & 176 & 160 & 150 & 43 & 110 & 142 & 64 & 249 & 192 & 70 & 17 & 90 & 83 \\
216 & 55 & 201 & 44 & 108 & 20 & 119 & 211 & 79 & 49 & 155 & 18 & 61 & 244 & 96 & 95 \\
14 & 234 & 152 & 247 & 113 & 168 & 146 & 233 & 94 & 220 & 123 & 242 & 191 & 76 & 253 & 250 \\
\end{tabular}\par}
\end{minipage}\par

\smallskip
For reproducibility, the following verifier accepts any of the moduli used here and the corresponding list in the order printed.
\begin{verbatim}
from math import gcd

def robust(n, a):
    units = {x for x in range(1, n) if gcd(x, n) == 1}
    vertices = set(a)
    edges = {(a[(i + 1) % len(a)] - x) % n for i, x in enumerate(a)}
    return (len(a) == len(vertices) == len(units) // 2
            and edges == vertices
            and (vertices | {-x % n for x in vertices}) == units)
\end{verbatim}
\end{proof}

To find these witnesses, choose a signed transversal $S$ with $\sum_{x\in S}x=0\pmod p$. This congruence is necessary because consecutive differences around a cycle sum to zero. Form three copies of $S$, for sources, targets and differences. For each $u,v\in S$ with $v-u\in S$, use the triple $(u,v,v-u)$ with one member in each copy. Choose triples covering every entry of every copy exactly once, discarding any choice that closes a directed cycle before all vertices belong to it. A completed choice gives a cyclic permutation with exactly the required difference set. The lists above allow independent verification without repeating this search.

\subsection{Three exceptional prime powers}
The list of unresolved robust constructions in Section~10 of Preece's chaplet paper includes $25$, $125$ and $289$. Here the underlying set is the unit group. A half-cycle on all nonzero residues of a composite modulus has a different ground set.

\begin{theorem}\label{thm:robust-three-powers}
There is a robust chaplet modulo $125$ and modulo $289$. There is none modulo $25$.
\end{theorem}

\begin{proof}
For the two positive cases, the lists below are cyclic. The vertex set $V$ equals the forward difference set $D$, and $V\cup(-V)=\U(n)$. The cardinalities of $V$, $D$ and $V\cup(-V)$ are respectively $(50,50,100)$ and $(136,136,272)$. Each displayed residue is a unit.

\par\noindent\begin{minipage}{\linewidth}
\noindent\textbf{Modulus $125$.}\par\smallskip
{\small\centering\setlength{\tabcolsep}{3pt}
\begin{tabular}{*{16}{r}}
124 & 17 & 79 & 91 & 98 & 106 & 32 & 6 & 67 & 76 & 24 & 56 & 62 & 14 & 16 & 83 \\
47 & 61 & 12 & 51 & 9 & 121 & 39 & 8 & 114 & 103 & 102 & 33 & 112 & 89 & 21 & 38 \\
71 & 18 & 77 & 43 & 59 & 37 & 84 & 57 & 28 & 72 & 96 & 99 & 2 & 73 & 94 & 7 \\
44 & 3 \\
\end{tabular}\par}
\end{minipage}\par

\par\noindent\begin{minipage}{\linewidth}
\noindent\textbf{Modulus $289$.}\par\smallskip
{\small\centering\setlength{\tabcolsep}{3pt}
\begin{tabular}{*{16}{r}}
288 & 115 & 63 & 162 & 45 & 206 & 120 & 150 & 215 & 285 & 21 & 239 & 105 & 48 & 36 & 38 \\
133 & 32 & 234 & 154 & 247 & 22 & 3 & 161 & 70 & 284 & 98 & 218 & 165 & 160 & 236 & 195 \\
66 & 176 & 103 & 212 & 185 & 108 & 270 & 201 & 211 & 193 & 167 & 198 & 143 & 188 & 250 & 109 \\
220 & 97 & 145 & 33 & 200 & 142 & 10 & 216 & 248 & 56 & 6 & 2 & 277 & 95 & 271 & 278 \\
126 & 280 & 166 & 157 & 197 & 260 & 276 & 209 & 231 & 47 & 8 & 90 & 61 & 229 & 175 & 183 \\
65 & 262 & 138 & 79 & 177 & 81 & 37 & 237 & 217 & 254 & 107 & 222 & 269 & 23 & 240 & 261 \\
110 & 243 & 99 & 25 & 168 & 155 & 263 & 235 & 7 & 30 & 111 & 275 & 171 & 93 & 172 & 137 \\
43 & 202 & 205 & 159 & 158 & 116 & 149 & 100 & 16 & 40 & 76 & 82 & 148 & 274 & 214 & 203 \\
64 & 164 & 24 & 62 & 245 & 230 & 31 & 232 \\
\end{tabular}\par}
\end{minipage}\par

\smallskip
The same verifier displayed above applies to both composite moduli.

Now suppose a robust chaplet modulo $25$ exists. Reversing its cyclic order interchanges $D$ and $-D$, so we can require $D=V$. A global negation, if necessary, ensures $1\in V$; it preserves $D=V$. Rotate the cycle to begin at $1$. Each of the other nine opposite pairs contributes one choice of sign, and the necessary condition $\sum_{v\in V}v\equiv0\pmod{25}$ leaves twenty possible sets. From $1$, extend a path only to an unused vertex of $V$ by an unused difference of $V$. Every possible normalized robust cycle would give a path of length ten. Across all twenty sets the numbers of paths with $1,\ldots,10$ vertices are
\[
       20,\ 70,\ 200,\ 448,\ 711,\ 805,\ 556,\ 238,\ 45,\ 0.
\]
In particular the last count vanishes even before testing the closing edge. The following complete program reproduces the counts.
\begin{verbatim}
from itertools import product

pairs = (2, 3, 4, 6, 7, 8, 9, 11, 12)
counts = [0] * 11

def visit(V, last, unused, seen):
    counts[10 - len(unused)] += 1
    for v in unused:
        d = (v - last) % 25
        if d in V and d not in seen:
            visit(V, v, unused - {v}, seen | {d})

for signs in product((1, -1), repeat=9):
    V = {1} | {(x*s) % 25 for x, s in zip(pairs, signs)}
    if sum(V) % 25 == 0:
        visit(V, 1, V - {1}, set())
assert counts[1:] == [20, 70, 200, 448, 711, 805, 556, 238, 45, 0]
\end{verbatim}
\end{proof}

\begin{corollary}\label{cor:robust}
Among odd $5<n<300$ that are not powers of three, the two and only two exceptions to existence of a robust chaplet are $17$ and $25$. For general odd $n>5$, a robust chaplet exists in each of the following cases:
\begin{enumerate}[label=(\alph*),nosep]
\item $3\mid n$ and $n$ has another prime divisor;
\item $n$ has a prime divisor $p\ge7$ with $p\equiv3\pmod4$;
\item $125$, $193$, $257$ or $289$ divides $n$.
\end{enumerate}
The odd positive integers admitting a robust chaplet have relative natural density one.
\end{corollary}
\begin{proof}
For primes in the stated range the assertion combines Preece's results as summarized in \cite[Section~6]{mo}, the nonexistence at $17$ proved there by exhaustive search, and Theorem~\ref{thm:robust-certs}. For the lifting claims, let $A$ and $D$ be the vertex and difference sets of a starting chaplet. If $D=\pm A$ for a fixed sign, taking products with the entire next unit group preserves this equality. Taking full inverse images under reduction also preserves it. Thus both operations in the proof of Theorem~\ref{thm:chaplet} preserve robustness.
\par
At the seed $3p$, the two cosets of the index-two subgroup $H$ are $H$ and $-H$. Therefore $(x-1)H$ equals one of them. This proves (a). For (b) with $3\nmid n$, the nonzero squares modulo $p$ form an odd-order subgroup $H$ of index two. The successive powers of a generator $x$ of $H$ give a cycle with $D=(x-1)H=\pm H$ (here $|H|\ge3$). Lift from this prime. For (c) start from the certified divisor whenever $3\nmid n$. The remaining cases are covered by (a). Here the composite seeds in Theorem~\ref{thm:robust-three-powers} may be used in exactly the same way as the prime seeds.
\par
To finish the range below $300$, an admissible $n$ divisible by $3$ is covered by (a). A remaining $n$ with any prime divisor other than $5$ and $17$ lifts from that robust prime. The possibilities left are $17,25,85,125,289$. Preece supplies $85$ in Note~5.2(a) of \cite{chaplet}. The two composite witnesses and the nonexistence proof at $25$ complete the classification, together with the published exclusion at $17$.
\par
For completeness, the reciprocals of the primes $p\equiv3\pmod4$ have divergent sum. To see this, take real $s>1$ and the nonprincipal character $\chi$ modulo four. Then
\[
 \log\zeta(s)-\log L(s,\chi)=2\sum_{p\equiv3\pmod4}p^{-s}+O(1)
 \qquad(s\downarrow1).
\]
The bounded error follows from the absolutely convergent prime-power terms of degree at least two and the single prime $2$. The left side tends to infinity because the alternating odd reciprocal series gives $L(1,\chi)=\pi/4$, with convergence as $s\downarrow1$ by the uniform alternating-series bound. Deleting $p=3$ preserves divergence. For any finite set of the remaining primes, CRT gives the relative density of odd integers divisible by none of them as $\prod_p(1-1/p)$. These products tend to zero. Part (b) proves the density assertion.
\end{proof}

\section{Keeping any fixed number of difference sequences}
The single change of multiplier also permits an explicit bound on the number of primitive residues excluded by higher differences. We include the bound because it still works when the order of the difference exceeds the length of the base cycle.
\begin{proposition}\label{thm:high}
Suppose $a,\Delta a,\ldots,\Delta^h a$ all list $\U(m)$ for an odd modulus $m>1$. If $p$ is an odd prime with $p\nmid m$ and $\varphi(p-1)>h(h+1)/2$, then for every $b\ge1$ there is such a word for $mp^b$.
\end{proposition}
\begin{proof}
Write $L=\varphi(m)$. For a generator $g$ modulo $p$, use as second coordinate at time $t\ge0$
\[
                    y_t=(-1)^t g^{\lfloor t/L\rfloor}\pmod p .
\]
This is periodic with period $L(p-1)$, and $y_{t+L}=g y_t$ since $L$ is even. For $0\le i<L$, the binomial formula for the $k$th forward difference gives
\[
 (\Delta^k y)_i=(-1)^{i+k}F_{i,k}(g),\qquad
 F_{i,k}(X)=\sum_{j=0}^k\binom{k}{j}X^{\lfloor(i+j)/L\rfloor}.
\]
In $\F_p[X]$, each polynomial is nonzero because $F_{i,k}(1)=2^k\ne0$. The sum of the degree bounds for a fixed $k$ is
\[
           \sum_{i=0}^{L-1}\lfloor(i+k)/L\rfloor=k,
\]
as one sees by division with remainder. Over all $1\le k\le h$, the union of forbidden roots has size at most $h(h+1)/2$. Thus one can choose a generator outside the union. Each $i$ now contributes the full multiplicative orbit of a nonzero second coordinate at every order $k\le h$, while the first coordinate is $(\Delta^k a)_i$. This proves the result for $mp$. Lemma~\ref{lem:repeated-prime-difference} increases the exponent to $b$.
\end{proof}
\section{An exact cycle count over finite rings}\label{sec:rings}
Let $R$ be a finite commutative ring of odd order greater than one, with identity. Write $V=R^\times$. Among permutations $T$ of $V$ for which $x\mapsto T(x)-x$ also permutes $V$, let $c(R)$ be the minimum number of cycles. Over a ring the answer can exceed one, even when every residue field has odd cardinality.
\begin{theorem}\label{thm:ring}
Let $s$ be the number of maximal ideals $\mathfrak m$ with $R/\mathfrak m\cong\F_3$. Then
\[
                        c(R)=2^{\max(s-1,0)}.
\]
\end{theorem}
\begin{proof}
Write $J$ for the intersection of the maximal ideals. They are pairwise comaximal, so CRT identifies $R/J$ with the product of all residue fields. An element is a unit precisely when it is in no maximal ideal. Thus $V$ is the full additive preimage of the product of their multiplicative groups.
\par
If $s\ge1$, project any admissible permutation edge onto the $s$ ternary coordinates. Its difference must be nonzero in each, so it goes from a sign vector $v$ to $-v$. A cycle stays inside a single pair $\{v,-v\}$. All $2^{s-1}$ such pairs occur in the image of $V$. For $s=0$ the lower bound of one is automatic.
\par
We first attain the bound in $R/J$. In the product of the ternary fields, take one two-cycle per sign pair. Each maps $v$ to $-v$ with difference $-2v=v$. The vertex and difference sets of each cycle are therefore the same pair. If $s=0$, start with a generator cycle on the units of any one residue field. All subsequent residue fields $K$ have at least five elements. Every existing cycle has even length. Attach $K^\times$ using the word of multipliers $-1,\ldots,-1,-g$ for a generator $g$ of $K^\times$. The product is $g$ and $-2,-g-1$ are nonzero in $K$. Lemma~\ref{lem:lift} applied separately on each cycle gives one lifted cycle over it. The vertex sets remain disjoint, as do their respective sets of labels, which together exhaust the units.
\par
It remains to lift through the additive group $J$. Refine it by a series of additive subgroups with quotients of prime order. At a single stage we have a surjective map of abelian groups with kernel $\langle z\rangle\cong\Z/p\Z$, $p$ odd. For an even base cycle $(x_i)_{i\bmod L}$ choose arbitrary lifts $X_i$ and put
\[
    T(X_i+t z)=X_{i+1}+(-t+b_i)z,\qquad
        b_0=\cdots=b_{L-2}=0,\quad b_{L-1}=1 .
\]
At the fibre of $X_0$, the return on $t$ is $t\mapsto t+1$, hence there is exactly one cycle upstairs. Its labels above $x_{i+1}-x_i$ are
\[
                   X_{i+1}-X_i+(-2t+b_i)z.
\]
They run through the entire fibre because $p$ is odd. Distinct base labels have disjoint fibres. The new length $pL$ is even, so iteration reaches $V$ while preserving the number of cycles and exhausting every difference exactly once. The additive subgroups used in this argument need no multiplicative structure.
\end{proof}
For example, $\F_3\times\F_3$ requires two cycles and $(\F_3)^3$ (the direct product of three copies) requires four. For $\Z/n\Z$ there is at most one residue field of order three, recovering Theorem~\ref{thm:sun}, including when the unit group has arbitrarily many generators.
\section{Sources, scope, and priority}
Preece explicitly asks his universal chaplet question and his separate prime-existence question in the closing paragraph of Section~10 \cite{chaplet}. Sun gives the odd-modulus conjecture after Theorem~1.4 \cite{sun}. Preece's earlier question about iterating differences introduced a further property and gave positive examples \cite[Section~6]{preece08}. Theorem~\ref{thm:two-orders} classifies the moduli for that property. Preece did not predict the exceptional value in that question.
\par
The closest explicit arithmetic constructions we located are Preece's prime-power and restricted composite daisy chains \cite{preece08}, his families from three generators \cite{preece09}, and the four-generator constructions of Preece and Vaughan \cite{pv}. The last paper also gives an example with five generators. The number of distinct prime divisors in our two main existence theorems is unrestricted. For instance, with $r$ distinct odd prime divisors, $\U(n)$ surjects onto $(\Z/2\Z)^r$ by the local quadratic characters. A construction requiring at most four cyclic factors cannot supply this family for $r>4$.
\par
The prime half-cycle of Buratti and Del Fra and the chaplet at $25$ are prior inputs. Preece's chaplet Theorems~3.6 and 3.7 already lift certain prime powers. Martin and Ollis study strong and robust half-cycles on full nonzero sets \cite{mo}. Robustness additionally identifies the vertex and difference sets up to an overall sign. Strength holds throughout Preece's admissible range. Robustness holds on a set of relative density one, and Corollary~\ref{cor:robust} classifies the entire admissible range below $300$.
\par
The cycle length over a base cycle is the classical net-voltage calculation \cite{gt}. Fibrewise assembly of orthomorphisms and prescribed cycle types also have precedents \cite{sw,bw}. In the corpus snapshot \cite{vmrc}, item 043 lifts a quotient matching by matching corresponding fibres, and item 057 decomposes unit groups into prime-power coordinates. These are the two construction ideas transferred here. The snapshot contains 225 numbered TeX manuscripts. Items 200 and 214 already contain the covering-radius and plactic-centralizer arguments examined during target selection. The arithmetic step here is the choice of the multipliers so that the return generates the fibre and every required first or second difference factor is invertible. The prime-power lift, the exceptional obstruction, and the higher-order polynomial bound specify the limits of this choice.
\par
\noindent\textbf{Assessment of the six criteria.}
G1 and G2 are witnessed by the original sources and the equivalence of sign conventions. For G3, the sources pose the general questions, and the later works examined establish restricted or related results. For G4, Theorems~\ref{thm:sun} and~\ref{thm:chaplet} prove the full statements without conjectural inputs. For G5, the constructions treat unbounded numbers of prime factors, classify second differences, and determine the precise cycle obstruction over odd rings. For G6, we checked the specialist papers above, their related terminology, and the corpus snapshot. This included follow-ups under the terms \emph{daisy chain}, \emph{chaplet}, \emph{robust half-cycle}, and \emph{orthomorphism}, as well as public GitHub issue and commit searches. We found no previous general solution to either universal question, or a previous robust witness at the two primes, in the material reviewed. That finding does not establish historical priority or settle the relative value of these results across different areas of mathematics.
\begin{thebibliography}{99}\small
\bibitem{preece08} D. A. Preece, Daisy chains---a fruitful combinatorial concept, \emph{Australas. J. Combin.} \textbf{41} (2008), 297--316.
\href{https://ajc.maths.uq.edu.au/pdf/41/ajc_v41_p297.pdf}{Article}.
\bibitem{sun-talk} Z.-W. Sun, Lecture at Zhejiang University, 15 November 2013, slide 34, conjecture dated 4 October 2013.
\href{https://www.mathnet.ru/PresentFiles/21090/sun.pdf}{Lecture slides}.
\bibitem{sun} Z.-W. Sun, Some new problems in additive combinatorics, \emph{Nanjing Univ. J. Math. Biquarterly} \textbf{36} (2019), 134--155, Theorem~1.4 and Remark~1.4.
\href{https://arxiv.org/abs/1309.1679}{arXiv:1309.1679}.
\bibitem{chaplet} D. A. Preece, Half-cycles and chaplets, \emph{Australas. J. Combin.} \textbf{43} (2009), 253--280.
\href{https://ajc.maths.uq.edu.au/pdf/43/ajc_v43_p253.pdf}{Article}.
\bibitem{preece09} D. A. Preece, Daisy chains with three generators, \emph{Australas. J. Combin.} \textbf{45} (2009), 157--174.
\href{https://ajc.maths.uq.edu.au/pdf/45/ajc_v45_p157.pdf}{Article}.
\bibitem{pv} D. A. Preece and E. R. Vaughan, Daisy chains with four generators, \emph{Australas. J. Combin.} \textbf{49} (2011), 77--93.
\href{https://ajc.maths.uq.edu.au/pdf/49/ajc_v49_p077.pdf}{Article}.
\bibitem{mo} G. N. Martin and M. A. Ollis, R-sequencings and strong half-cycles from narcissistic terraces, \emph{Australas. J. Combin.} \textbf{63} (2015), 346--367.
\href{https://ajc.maths.uq.edu.au/pdf/63/ajc_v63_p346.pdf}{Article}.
\bibitem{gt} J. L. Gross and T. W. Tucker, \emph{Topological Graph Theory}, Wiley, 1987, Section~2.1.3.
\bibitem{sw} D. S. Stones and I. M. Wanless, Compound orthomorphisms of the cyclic group, \emph{Finite Fields Appl.} \textbf{16} (2010), 277--289.
\href{https://doi.org/10.1016/j.ffa.2010.04.001}{DOI}.
\bibitem{bw} A. Bors and Q. Wang, Coset-wise affine functions and cycle types of complete mappings, 2021.
\href{https://arxiv.org/abs/2109.03922}{arXiv:2109.03922}.
\bibitem{vmrc} QEEDDD, \emph{Verifiable Mathematical Reasoning Corpus}, items 043 and 057, snapshot of 20 November 2026, commit \texttt{c2ac34ea4622}.
\href{https://github.com/QEEDDD/Verifiable-Mathematical-Reasoning-Corpus/tree/c2ac34ea4622778187173f41bb1996deee8081b7}{Pinned repository}.
\end{thebibliography}
\end{document}
